Finding amino acid sequences based on DNA sequences is one of the most
basic procedures in bioinformatics. The program \texttt{blastx} is a
golden standard for this task. Users usually filter and format the
output of \texttt{blastx} using custom scripts according to their
research goals. One such goal is to find peptide sequences given
reference protein sequences and some query DNA regions. The
program \texttt{getpep} streamlines the basic usage of \texttt{blastx}
for this purpose.

The program \texttt{getpep} accepts a \texttt{fasta} file of
nucleotide sequences and a \texttt{fasta} file of protein
sequences. It builds a command for \texttt{blastx}, runs it, filters
the output for the best hits per query, applies the user-specified
identity threshold, and returns the hits in one of the three formats:

\begin{enumerate}
  \itemsep0em
  \item whole \texttt{fasta} sequences of proteins, for which hits were found;
  \item hit \texttt{fasta} sequences of peptides (with coordinates);
  \item a filtered output of \texttt{blastx} hits.
\end{enumerate}
